\section{Statements}

\subsection{Ethics \& Broader Impacts}
\label{appx:statement:impacts}
Prior works have established how the growing usage of LMs that can code also carries a certain amount of risk.
We identify three main risks that could arise when building and using a system like SWE-agent, then discuss how we incorporates measures that mitigate such problems.

First is the security risks that come with executing LM-generated code on device.
When evaluating on software engineering tasks that feature an incredibly diverse number of issue descriptions, running code generations on a personal computing device can have negative side effects, such as the unintentional removal of digital assets (e.g., \texttt{rm -rf asset/}).
To defend against this, we design SWE-agent to use ephemeral containers for both inference and evaluation.
SWE-agent's execution environment and the SWE-bench evaluation framework are both carried out in sand-boxed code environments, which is made possible with Docker.
Executing code in a Docker container ensures that its effects are mostly isolated from the rest of the system.
While not considered as secure as virtualized hardware isolation, the namespace isolation provided by Docker containers is deemed sufficient for code that is not deliberately engineered to exploit recent container escape vulnerabilities.
More details are discussion is in \S\ref{appx:interface:implementation}.

Second, if the wider community develops interest for SWE-agent and builds upon it, it is also possible that illegitimate evaluation datasets or infrastructure can be used to inject testing devices with malicious code or instructions to generate malicious code.
For instance, an unofficial repository claiming to host an inference/evaluation harness for SWE-agent/bench could include a task instance with an issue description that tells the LM agent to build key logging functionality and store it in a hidden folder.
To eliminate confusion and reduce the possibility of such an event, we provide clear guidelines listed on our GitHub repositories, data stores, and websites indicating the official repositories and channels that we actively maintain.
We also encourage third parties to incorporate any improvements into our codebase and help with integrating such contributions.

Lastly are the consequences of software engineering agents being deployed in the real world.
Prior works have conceptualized and put forth prototypes of agents that can carry out offensive security measures.
It is also not difficult to imagine that a system like SWE-agent can be incorporated into pipelines resulting in the production of malicious code.
SWE-agent’s strong performance on SWE-bench implies that future AI systems will likely be increasingly adept in the aforementioned use cases.
Releasing SWE-agent as an open source tool can support research towards designing sound, effective constraints for what software engineering agents are permitted to do.
It can also serve as a system that legal experts and policy-making entities can experiment with to shape the future of what AI-driven end to end software engineering could look like.

\subsection{Reproducibility}
\label{appx:statement:reproducibility}
To help the greater community reproduce the results presented in this paper and build on the SWE-agent platform, we open source all of our resources that were created for this project.
The source code for the interactive pipeline, context management logic, command implementations, interface design, and everything else is entirely available in a GitHub repository.
We provide extensive text and video documentation describing how to run and modify different parts of the codebase.
Practitioners should be able to easily recover our findings by running the agent with simple scripts.
We also open source all inference and evaluation artifacts (e.g., trajectories, code generations, evaluation execution traces, analysis notebooks).
The results presented in the main and supplementary parts of this paper can be fully rendered from the data.
Finally, we also maintain an active online help forum to assist with any reproduction problems or questions about how to build on ACI design and SWE-agent.

\subsection{Limitations \& Future Work}
\label{appx:statement:limits}
The final SWE-agent configuration has a small toolkit, albeit highly effective.
With SWE-agent's highly extensible design, we're excited by the prospect of adding more tools, such as web browsing or static analysis, that can leverage more signals from an issue description and codebase to improve the \% Resolved performance.
Many tools trialed by prior works from software engineering and language model agents, such as static/dynamic analysis, spectrum based fault localization, or test generation via fuzzing could prove useful.

Second, in this work, the ACI development process and case studies are done manually.
Many components of SWE-agent were crafted from observations of recurring behavior within a single trajectory or across multiple trajectories.
Automating part or all of this process could not only accelerate work built on top of SWE-agent, but also provide greater insights into developing ACI principles for agentic software engineering.
Contemporary works have explored automated prompting to improve performance on traditional sequence to sequence tasks, supplanting the need for manual prompt design.
Thinking about automating ACI design raises immediately interesting questions around how such systems can scrutinize and iterate upon their own designs. Ensuring such horizon leads to incremental performance improvements across a longer horizon is also a challenging question.

Finally, the scope of SWE-agent is exclusively focused on programmatic tasks like software engineering and code generation.
We're curious to see whether the same principles of ACI and our observations of agent behavior are transferable to different domains.
Recent work around applying LM agents to a variety of digital work applications have proliferated, such as use cases in education technology, data analysis, and enterprise workflows.
We hope that thinking about improving performance of agentic workflows on these domains through the lens of ACI design can be a symbiotic process.
For instance, for a task such a shopping on the web, in place of a typical Google-style search tool, could agents benefit from additional information beyond a list of each page’s title and snippet?
Would the design vary if the nature of the downstream task were to change slightly?
For a completely different task, such as navigating an internal company knowledge base to help a recently on-boarded employee, how might the search interface be best adjusted to the agent?

Similar to the progression of the field of User Experience (UX) and Human Computer Interaction (HCI) research, applying ACI to other domains could not only yield improvements in downstream task performance, but also further expand the list of ACI principles.
We believe that the fundamental motivations for ACI, the foundational principles we put forth, and our case study of SWE-agent as an instantiation of implementing and improving an ACI can motivate such work.